\section{Introduction}

Automated theorem proving (ATP) for first-order logic is the problem of
mechanically determining whether a formula follows from given axioms.
Industrial-strength provers such as Vampire~\cite{kovacs2013vampire} and
E~\cite{schulz2002E} use saturation-based resolution and superposition
calculi; the connection-based leanCoP~\cite{otten2003leancop} avoids
explicit instantiation through matrix methods and unification. These
systems solve thousands of problems from the TPTP
library~\cite{Sut17} and are applied in hardware verification, software
model-checking, and proof-assistant automation. The pedagogical LK$'$
calculus of Hou~\cite[Fig.~2.3, p.~65]{hou2021fundamentals} occupies a
different position: it is a backward sequent search whose rules
can be stated in twenty lines of pseudocode, making it suitable for formal
study, but whose na\"ive rule-application order can lead to pathological
search behaviour. LK$'$ eliminates the structural rules of
LK~\cite{gentzen1935untersuchungen} and makes $\forall L$ and $\exists R$
\emph{reusable}: the quantified formula is retained after
instantiation so the same axiom can be applied with multiple witnesses on a
single branch, and the non-deterministic choice of witness is the principal
source of search-space explosion.

This paper identifies two structural failure modes in the na\"ive LK$'$
backward search of Hou~\cite[Algorithm~2, p.~67]{hou2021fundamentals} and
proposes two targeted improvements: a six-class \emph{priority schedule}
that refines rule-application order (Sect.~\ref{sec:approach}) and an
\emph{iterative-deepening} envelope around the depth-first
search~\cite{korf1985iterative}. An empirical evaluation on three datasets
(888~TPTP problems, 280~FOLIO problems, 202~synthetic) shows consistent
improvement in decided problems across all problem families, with gains
concentrated where short proofs exist but are blocked by deep
quantifier-instantiation chains (Sect.~\ref{sec:experiments}).
